% GENERATED FILE. Edit canonical lessons, then run npm run generate:books.
% canonical-source-hash: fnv1a64:fc00e40bf0073971
% canonical-lessons: MW-C33-hear-chha, MW-C33-chha, MW-C33-hear-saat, MW-C33-saat, MW-C33-hear-aath, MW-C33-aath, MW-C33-hear-no-nine, MW-C33-no-nine, MW-C33-hear-das, MW-C33-das, MW-R33-count-recall, MW-C33-count-ten, MW-R33-script-close

\chapter{Six To Ten}
\label{ch:mw-count-ten}
\hlchaptermodality{33}

\begin{chapteropening}
\textbf{By the end of this chapter:} \emph{I can hear, say, read, and write every number from one to ten, in any order.}

Five more numbers arrive with no new sign at all, and the whole run of ten is then scored shuffled in each of the four skills.
\end{chapteropening}

\section[chha]{Hear \emph{chha} — six}

\label{lesson:MW-C33-hear-chha}



\begin{quote}
Recall the counting payoff, say tea, then say four.

\begin{itemize}
\raggedright
  \item \emph{Recall:} \emph{chār} from the last chapter, and the plain \emph{ch} you tested with a hand at your mouth
\end{itemize}
\end{quote}

\subsection*{What to know first}
\begin{quote}
\emph{chha} — \textbf{six}
\end{quote}

Hand at the mouth again. \emph{chār} moves nothing. \emph{chha} pushes a puff of air out with the consonant.

That single difference is the whole of it. The tongue does the same thing in both words; the lungs do not. Marwari keeps that contrast everywhere, and here it separates two numbers that sit next to each other on a price ticket.

The same number is also printed \emph{chhe}, with a vowel mark after the consonant. Both are in use; count with the bare form for now and read the other when a page prints it.

\begin{tcolorbox}[breakable,colback=teal!4,colframe=teal!35!black,title={Before you move on}]
Sources: \href{https://www.omniglot.com/language/numbers/rajasthani.htm}{Omniglot, Numbers in Rajasthani}; \href{https://dsal.uchicago.edu/dictionaries/soas/}{Turner, A Comparative Dictionary of the Indo-Aryan Languages}.
\end{tcolorbox}
\section[chha]{\mw{छ} — six on the page}

\label{lesson:MW-C33-chha}



\begin{quote}
Write \textbf{\mw{य}}, write \textbf{\mw{ए}}, then write four.
\end{quote}

\subsection*{What to know first}
\begin{quote}
\textbf{\mw{छ}} — \emph{chha} — six
\end{quote}

One sign, made in chapter 9, and nothing added to it.

The other printed spelling is \textbf{\mw{छे}}, and it is worth a moment. That is \textbf{\mw{छ}} with the mark of the letter chapter 32 taught — the same vowel, riding on a consonant instead of standing alone. It is also, exactly, the second syllable of \textbf{\mw{पाछे}}. A chunk the hand has been writing since chapter 9 turns out to be a number.

\subsection*{Your turn}
Write \textbf{\mw{छ}}, then write \textbf{\mw{छे}} beneath it, then write \textbf{\mw{ए}} beside both. Say which of the three stands alone and which two need something under them.

\begin{tcolorbox}[breakable,colback=teal!4,colframe=teal!35!black,title={Before you move on}]
Sources: \href{https://www.omniglot.com/language/numbers/rajasthani.htm}{Omniglot, Numbers in Rajasthani}; \href{https://www.unicode.org/charts/PDF/U0900.pdf}{Unicode Devanagari chart}.
\end{tcolorbox}
\section[sāt]{Hear \emph{sāt} — seven}

\label{lesson:MW-C33-hear-saat}



\begin{quote}
Give the greeting this book opened with, then say six and three.

\begin{itemize}
\raggedright
  \item \emph{Recall:} \emph{tīn} and its long vowel, two chapters back
\end{itemize}
\end{quote}

\subsection*{What to know first}
\begin{quote}
\emph{sāt} — \textbf{seven}
\end{quote}

The first syllable is \emph{sā}. That is the respect particle at the end of \textbf{\mw{राम}-\mw{राम} \mw{सा}}, the first thing this book ever taught. The sound has been in the ear since page one; only its job here is new.

Seven closes on the dental \emph{t} of \emph{tīn}, not the retroflex one. Front of the mouth at both ends.

\begin{tcolorbox}[breakable,colback=teal!4,colframe=teal!35!black,title={Before you move on}]
Sources: \href{https://www.omniglot.com/language/numbers/rajasthani.htm}{Omniglot, Numbers in Rajasthani}; \href{https://dsal.uchicago.edu/dictionaries/soas/}{Turner, A Comparative Dictionary of the Indo-Aryan Languages}.
\end{tcolorbox}
\section[sāt]{\mw{सात} — seven on the page}

\label{lesson:MW-C33-saat}



\begin{quote}
Write tea, then write six.
\end{quote}

\subsection*{What to know first}
\begin{quote}
\textbf{\mw{सा}} + \textbf{\mw{त}} $\to$ \textbf{\mw{सात}} — \emph{sāt} — seven
\end{quote}

\textbf{\mw{सा}} was the second word-piece this book asked the hand to build, in chapter

\begin{enumerate}
\raggedright
  \item Seven is that piece with a dental \textbf{\mw{त}} after it, and nothing else.
\end{enumerate}

\subsection*{Your turn}
Look at \textbf{\mw{सात}}, cover it, wait five seconds, write it. Then write \textbf{\mw{तीन}} under it and mark the \textbf{\mw{त}} in both.

\begin{tcolorbox}[breakable,colback=teal!4,colframe=teal!35!black,title={Before you move on}]
Source: \href{https://www.omniglot.com/language/numbers/rajasthani.htm}{Omniglot, Numbers in Rajasthani}.
\end{tcolorbox}
\section[āṭh]{Hear \emph{āṭh} — eight}

\label{lesson:MW-C33-hear-aath}



\begin{quote}
Recall the seven-word food payoff, offer formal thanks, then say seven.

\begin{itemize}
\raggedright
  \item \emph{Recall:} \emph{chār} and \emph{pāṅch}, and which of the two carries a nasal
\end{itemize}
\end{quote}

\subsection*{What to know first}
\begin{quote}
\emph{āṭh} — \textbf{eight}
\end{quote}

It opens on the long \emph{ā} of \textbf{\mw{आभार}}, standing on its own with no consonant in front of it. Of the ten numbers this slice teaches, eight is the only one that does that.

The closing consonant is retroflex: the tongue curls back, as in \textbf{\mw{ठंडी}}. Do not close it at the teeth, or seven and eight begin to sound alike.

\begin{tcolorbox}[breakable,colback=teal!4,colframe=teal!35!black,title={Before you move on}]
Sources: \href{https://www.omniglot.com/language/numbers/rajasthani.htm}{Omniglot, Numbers in Rajasthani}; \href{https://dsal.uchicago.edu/dictionaries/soas/}{Turner, A Comparative Dictionary of the Indo-Aryan Languages}.
\end{tcolorbox}
\section[āṭh]{\mw{आठ} — eight on the page}

\label{lesson:MW-C33-aath}



\begin{quote}
Write formal thanks, then write seven.
\end{quote}

\subsection*{What to know first}
\begin{quote}
\textbf{\mw{आ}} + \textbf{\mw{ठ}} $\to$ \textbf{\mw{आठ}} — \emph{āṭh} — eight
\end{quote}

Two letters, both of them years old in this book's own timeline: \textbf{\mw{आ}} from chapter 2 and \textbf{\mw{ठ}} from chapter 17. Nothing is joined, nothing is marked.

Compare \textbf{\mw{सात}} and \textbf{\mw{आठ}} side by side. Seven takes the vowel as a mark on a consonant; eight takes the same vowel as a letter, because there is no consonant for it to sit on.

\subsection*{Your turn}
Look at \textbf{\mw{आठ}}, cover it, wait five seconds, write it. Then write \textbf{\mw{सात}} above it and circle the long \emph{ā} in each.

\begin{tcolorbox}[breakable,colback=teal!4,colframe=teal!35!black,title={Before you move on}]
Source: \href{https://www.omniglot.com/language/numbers/rajasthani.htm}{Omniglot, Numbers in Rajasthani}.
\end{tcolorbox}
\section[no]{Hear \emph{no} — nine}

\label{lesson:MW-C33-hear-no-nine}



\begin{quote}
Say this one, say very, then say eight.

\begin{itemize}
\raggedright
  \item \emph{Recall:} \emph{ek} and the letter \textbf{\mw{ए}} you had to learn to write it
\end{itemize}
\end{quote}

\subsection*{What to know first}
\begin{quote}
\emph{no} — \textbf{nine}
\end{quote}

One consonant, one long vowel, and it rhymes with \emph{do}. Two and nine are the only numbers in this run that short.

This is not the English word \textbf{no}. Marwadi has not yet given this book any word for refusing at all, which is a real gap and not a spelling accident: a learner who mishears \emph{no} as a refusal will walk away from a fair price.

\begin{tcolorbox}[breakable,colback=teal!4,colframe=teal!35!black,title={Before you move on}]
Sources: \href{https://www.omniglot.com/language/numbers/rajasthani.htm}{Omniglot, Numbers in Rajasthani}; \href{https://dsal.uchicago.edu/dictionaries/soas/}{Turner, A Comparative Dictionary of the Indo-Aryan Languages}.
\end{tcolorbox}
\section[no]{\mw{नो} — nine on the page}

\label{lesson:MW-C33-no-nine}



\begin{quote}
Write this one, then write eight.
\end{quote}

\subsection*{What to know first}
\begin{quote}
\textbf{\mw{न}} + \textbf{\mw{ो}} $\to$ \textbf{\mw{नो}} — \emph{no} — nine
\end{quote}

Write \textbf{\mw{दो}} and \textbf{\mw{नो}} one under the other. The mark on the right is the same in both. Only the letter carrying it changes, and with it the number.

\subsection*{Your turn}
Look at \textbf{\mw{नो}}, cover it, wait five seconds, write it. Then write \textbf{\mw{दो}} beside it and read the pair aloud without looking.

\begin{tcolorbox}[breakable,colback=teal!4,colframe=teal!35!black,title={Before you move on}]
Source: \href{https://www.omniglot.com/language/numbers/rajasthani.htm}{Omniglot, Numbers in Rajasthani}.
\end{tcolorbox}
\section[das]{Hear \emph{das} — ten}

\label{lesson:MW-C33-hear-das}



\begin{quote}
Say show it to me, say five, then say nine.

\begin{itemize}
\raggedright
  \item \emph{Recall:} \emph{chha} and the puff of air that separates it from \emph{chār}
\end{itemize}
\end{quote}

\subsection*{What to know first}
\begin{quote}
\emph{das} — \textbf{ten}
\end{quote}

Short vowel this time, closed off at once by \emph{s}. Nothing is held.

\emph{das} is the first number in this book that a Rajasthani market is likely to say out loud. Almost nothing on a stall is priced under ten, which is why a course that stopped at five could ask the price four different ways and still not understand the answer.

The comparative dictionary of the Indo-Aryan languages lists this form for Marwari by name, under the Sanskrit word for ten.

\begin{tcolorbox}[breakable,colback=teal!4,colframe=teal!35!black,title={Before you move on}]
Sources: \href{https://www.omniglot.com/language/numbers/rajasthani.htm}{Omniglot, Numbers in Rajasthani}; \href{https://dsal.uchicago.edu/dictionaries/soas/}{Turner, A Comparative Dictionary of the Indo-Aryan Languages}.
\end{tcolorbox}
\section[das]{\mw{दस} — ten on the page}

\label{lesson:MW-C33-das}



\begin{quote}
Write \textbf{\mw{ख}}, then write nine.
\end{quote}

\subsection*{What to know first}
\begin{quote}
\textbf{\mw{द}} + \textbf{\mw{स}} $\to$ \textbf{\mw{दस}} — \emph{das} — ten
\end{quote}

Bare consonants. No mark rides above the line and none hangs below it, because the vowel between them is the short \emph{a} that every Devanagari consonant carries without being told.

The \textbf{\mw{द}} is two's; the \textbf{\mw{स}} is seven's. Ten is written out of pieces the hand already made for other numbers.

\subsection*{Your turn}
Look at \textbf{\mw{दस}}, cover it, wait five seconds, write it. Then write \textbf{\mw{दो}} beside it and say what the mark on the second one does to the vowel.

\begin{tcolorbox}[breakable,colback=teal!4,colframe=teal!35!black,title={Before you move on}]
Source: \href{https://www.omniglot.com/language/numbers/rajasthani.htm}{Omniglot, Numbers in Rajasthani}.
\end{tcolorbox}
\section[Practice]{Ten words, shuffled}

\label{lesson:MW-R33-count-recall}



\begin{quote}
Write the show-request word, recall the counting payoff from the last chapter, then write \textbf{\mw{ए}}.
\end{quote}

\subsection*{Your turn}
Two runs of five are still two runs until they are mixed.

\begin{enumerate}
\raggedright
  \item Hear ten numbers named singly, drawn from across the whole set, and say each.
  \item Count backwards from ten to one.
  \item Read \textbf{\mw{आठ}}, \textbf{\mw{तीन}}, \textbf{\mw{दस}}, \textbf{\mw{छ}}, \textbf{\mw{एक}}, \textbf{\mw{नो}} printed in that order and say each.
  \item Write \textbf{\mw{सात}}, \textbf{\mw{पांच}}, \textbf{\mw{नो}}, \textbf{\mw{दो}}, \textbf{\mw{दस}} from dictation.
\end{enumerate}

Then say which pairs share a sign: two and nine, seven and ten, four and five.

\begin{tcolorbox}[breakable,colback=teal!4,colframe=teal!35!black,title={Before you move on}]
Source: \href{https://www.omniglot.com/language/numbers/rajasthani.htm}{Omniglot, Numbers in Rajasthani}.
\end{tcolorbox}
\section[Practice]{Payoff — counting to ten}

\label{lesson:MW-C33-count-ten}



\begin{quote}
Recall the five-number payoff, then write the letter that made it possible.
\end{quote}

\subsection*{Your turn}
Four skills, scored separately.

\begin{enumerate}
\raggedright
  \item \textbf{Listening.} Hear ten numbers named singly and out of order; write each in figures of your own language.
  \item \textbf{Speaking.} Count one to ten, then ten to one, with no model.
  \item \textbf{Reading.} Read all ten printed in Devanagari, shuffled, and say each.
  \item \textbf{Writing.} Write all ten from dictation, shuffled, with no model in view.
\end{enumerate}

Pass each one separately.

Look back at what the hand had to learn for this: one letter, \textbf{\mw{ए}}. Every other sign in all ten words was already in the book. A reader who felt the numbers were going to be a wall of new writing has the evidence in front of them that they were not.

\begin{tcolorbox}[breakable,colback=teal!4,colframe=teal!35!black,title={Before you move on}]
Sources: \href{https://www.omniglot.com/language/numbers/rajasthani.htm}{Omniglot, Numbers in Rajasthani}; \href{https://www.marwaripathshala.com/marwari-lesson-9-english}{Marwari Pathshala, Lesson 9}.
\end{tcolorbox}
\section[Practice]{Close the writing loop for \mw{छ}, \mw{ठ} and \mw{स}}

\label{lesson:MW-R33-script-close}



\begin{quote}
Recall the ten-number payoff, say the show request, then write market.
\end{quote}

\subsection*{Your turn}
Write three single signs from sound cues: \textbf{\mw{छ}}, \textbf{\mw{ठ}}, \textbf{\mw{स}}. Then write \textbf{\mw{च}} beside \textbf{\mw{छ}} and mark which of the two carries the breath.

Then write \textbf{\mw{छ}}, \textbf{\mw{सात}}, \textbf{\mw{आठ}}, \textbf{\mw{नो}}, \textbf{\mw{दस}} from dictation. Repair only the missed sign and rewrite its word once.

\begin{tcolorbox}[breakable,colback=teal!4,colframe=teal!35!black,title={Before you move on}]
Sources: \href{https://www.omniglot.com/language/numbers/rajasthani.htm}{Omniglot, Numbers in Rajasthani}; \href{https://www.unicode.org/charts/PDF/U0900.pdf}{Unicode Devanagari chart}.
\end{tcolorbox}
